\begin{example}
   Let $n = 10$. We have $4 + 9 \equiv 3 \pmod{10}$, $-5 \equiv 5 \pmod{10}$ and $4 - 7 \equiv 7 \pmod{10}$. With multiplication, $4 \cdot 9 \equiv 6 \pmod{10}$ and $4 \cdot (-9) \equiv 4 \pmod{10}$. Remark that the result of the operation is always the remainder of the operation when divided by $10$. 
\end{example}
\begin{remark}
   The result of any modular operation is its remainder when divided by $n$. 
\end{remark}

\subsection{Solving Linear Equations}

Naturally, the question arises, how do we solve equations of the form
\[ ax \equiv b \pmod{n} \]
if a solution even exists? For example, let $n = 10$, does $9x \equiv 1$ have a solution? Yes, indeed, we can take $x = 9$, since $9 \cdot 9 = 81 \equiv 1 \pmod{10}$. What about $4x \equiv 1$? This equation does not have a solution due to the following proposition:

\begin{propbox}
   Let $a$, $b \in \ZZ$, and $n \in \NN^+$. Let $d := \gcd(a,n)$. The equation $ax + b \equiv 0 \pmod{n}$ has a solution if and only if $d \mid b$. 
\end{propbox}
\begin{proof}
   Let $ax + b \equiv 0 \pmod{n}$ have a solution, then $b = kn - ax$ for some $k \in \ZZ$, and so $d \mid b$. Conversely, let $d \mid b$ and consider the equation $(a/d)x + (b/d) \equiv 0 \pmod{n/d}$. A solution to this equation also solves the original equation $ax + b \equiv 0 \pmod{n}$. We have that $\gcd(a/d, n/d) = 1$, so the equation has a unique solution. 
\end{proof}

\begin{remark}
   If $d > 1$, then there are multiple solutions to the linear equation $ax + b \equiv 0 \pmod{n}$. 
\end{remark}

\begin{example}
   Let $n = 10$. The equation $4x \equiv 1$ does not have a solution since $\gcd(4, 10) = 2$ and $2 \nmid 1$. However, $9x \equiv 1$ has a solution since $\gcd(9,10) = 1$ and $1 \mid 1$. 
\end{example}

We know that $ax \equiv b \pmod{n}$ has exactly one solution when $\gcd(a,n) = 1$. In this case, how do we find $\ol{a^{-1}} \in \ZZ/n\ZZ$? We want to find an $u \in \ZZ$ such that $u \cdot a \equiv 1 \pmod{n}$. To do so, we can use Theorem 6.1.2 to write:
\[ 1 = ua + vb \]
for $u$, $v \in \ZZ$. Since $ua + vb \equiv ua \pmod{n}$, we have that $\ol{u} = \ol{a^{-1}}$. 

\begin{example}
   We want to find $\ol{4^{-1}}$ in $\ZZ/15\ZZ$.
\end{example}
\begin{solution}
   We know a solution exists because $\gcd(4, 15) = 1$. By applying the Euclidean algorithm and back substitution, we obtain:
   \[ 1 = 4 \cdot 4 + (-1) \cdot 15 \equiv 4 \cdot 4 \pmod{15} \]
   which implies that $\ol{4^{-1}} = \ol{4}$. 
\end{solution}

\subsection{Inverses and Zero Divisors}

\begin{defbox}
   Let $\ol{a} \in \ZZ/n\ZZ$.
   \begin{enumerate}
      \item Call $\ol{b} \in \ZZ/n\ZZ$ a \textbf{multiplicative inverse of $\ol{a}$} if $\ol{a} \cdot \ol{b} = \ol{1} = \ol{b} \cdot \ol{a}$. If $\ol{a}$ has an inverse, we say $\ol{a}$ is \textbf{invertible}. 
      \item Call $\ol{a}$ a \textbf{zero divisor} if $\ol{a} \neq \ol{0}$ and there exists a $\ol{b} \in \ZZ/n\ZZ$, $\ol{b} \neq \ol{0}$, such that $\ol{a} \cdot \ol{b} = \ol{0}$ or $\ol{b} \cdot \ol{a} = \ol{0}$. 
   \end{enumerate}
\end{defbox}
\begin{remark}
   An inverse, if it exists, is unique. Indeed, if $\ol{b_0}$ and $\ol{b_1}$ are inverses of some $\ol{a} \in \ZZ/n\ZZ$, then:
   \[ \ol{b_0} = \ol{b_0} \cdot \ol{1} = \ol{b_0} \cdot \ol{(a \cdot b_1)} = \ol{(b_0 \cdot a)} \cdot \ol{b_1} = \ol{1} \cdot \ol{b_1} = \ol{b_1}. \]
   Because of this uniqueness, we usually denote the inverse of $\ol{a}$ as $\ol{a^{-1}}$. 
\end{remark}

\begin{propbox}
   Let $n \geq 2$ be a natural number and $\ol{a} \in \ZZ/n\ZZ$ such that $\ol{a} \neq \ol{0}$. $\ol{a}$ is invertible if and only if $\ol{a}$ is not a zero divisor. 
\end{propbox}
\begin{proof}
   Let $\ol{a}$ be invertible. It suffices to show that for every $\ol{b} \in \ZZ/n\ZZ$, if $\ol{a} \cdot \ol{b} = 0$ or $\ol{b} \cdot \ol{a} = 0$, then $\ol{b} = 0$. Suppose $\ol{a} \cdot \ol{b} = 0$. Then:  
   \[ \ol{b} = \ol{1} \cdot \ol{b} = (\ol{a^{-1}} \cdot \ol{a}) \cdot \ol{b} = \ol{a^{-1}} \cdot (\ol{a} \cdot \ol{b}) = \ol{a^{-1}} \cdot 0 = 0. \]
   The case where $\ol{b} \cdot \ol{a} = 0$ is proven analogously. Conversely, suppose $\ol{a}$ is not a zero divisor. Consider the set $S := \{\ol{a} \cdot \ol{1}, \ol{a} \cdot \ol{2}, \dots, \ol{a} \cdot \ol{(n-1)}$. We know that $\ol{0} \notin S$. Suppose by contradiction that $\ol{1} \notin S$, then by the pigeonhole principle, there exists $i$ and $j$ such that $1 \leq i < j \leq n - 1$ and $\ol{a} \cdot \ol{i} = \ol{a} \cdot \ol{j}$. But then $\ol{a} \cdot \ol{(j - i)} = \ol{0}$ and $\ol{j - i} \neq \ol{0}$, so $\ol{a}$ is a zero divisor, contradiction. 
\end{proof}
